\documentclass[
	fontsize=12pt,          % default font size 12pt
	paper=a4,               % DIN A4 page format
	numbers=noenddot,       % remove dots behind chapter numbers (e.g. 1.5 not 1.5.)
	listof=totoc,           % add list of figures, tables, etc. to ToC
	listof=entryprefix,     % add entry name to figures, tables, etc.
	listof=nochaptergap,    % no chapter gap for figures, tables, etc.
	bibliography=totoc,     % add bibliography to ToC but without a chapter number
	parskip=half            % half line spacing between paragraphs
	openany                 % chapters can start on any page
]{scrbook}
\usepackage{tcolorbox}

\include{../../for_latex/preamble}
\title{\Huge\textbf{Auflagerkräfte}}
\author{Luis Staudt}
\date{}

% Code-Highlighting für Java
\definecolor{codegray}{rgb}{0.5,0.5,0.5}
\definecolor{codepurple}{rgb}{0.58,0,0.82}
\definecolor{backcolour}{rgb}{0.95,0.95,0.92}
\definecolor{keywordcolor}{rgb}{0.8,0.47,0.13}

\lstdefinestyle{java}{
	language=,
	basicstyle=\ttfamily\small,
	keywordstyle=\bfseries\color{blue!70!black},
	stringstyle=\color{red!60!black},
	commentstyle=\itshape\color{gray},
	numbers=left,
	numberstyle=\tiny\color{gray},
	numbersep=8pt,
	frame=single,
	framerule=0.4pt,
	rulecolor=\color{gray!50},
	breaklines=true,
	showstringspaces=false,
	tabsize=4,
	xleftmargin=1.5em,
	framexleftmargin=1.5em,
	aboveskip=0.8em,
	belowskip=0.6em
}
\lstset{style=java}

\newtcolorbox{hinweis}[1][Hinweis]{
	colback=blue!5,
	colframe=blue!40!black,
	fonttitle=\bfseries,
	title=#1,
	boxrule=0.5pt,
	arc=2pt
}

\newtcolorbox{beispiel}[1][Beispiel]{
	colback=green!5,
	colframe=green!40!black,
	fonttitle=\bfseries,
	title=#1,
	boxrule=0.5pt,
	arc=2pt
}

\begin{document}

% --- Titelblock ---
\begin{center}
	\rule{\textwidth}{0.8pt}\\[0.6em]
	{\LARGE\bfseries Intensivtutorium}\\[0.3em]
	{\large Programmierung}\\[0.6em]
	\begin{tabular}{r l @{\hspace{2cm}} r l}
		\textbf{Semester:}      & Sommersemester 2026                    & \textbf{Tutor:}         & Luis Staudt \\
		\textbf{Studiengänge:}  & MIB, OMB, PEB, MVB         & & \\
	\end{tabular}\\[0.6em]
	\rule{\textwidth}{0.8pt}
\end{center}
\vspace{1em}

% --- Seitenkopf und -fuß (ab Seite 2) ---
\pagestyle{fancy}
\fancyhf{}
\lhead{\small Intensivtutorium -- Grundlagen der Programmierung}
\rhead{\small SS 2026}
\cfoot{\thepage}
\renewcommand{\headrulewidth}{0.4pt}


% =============================================
% Übungsaufgabe -- Auflagerkräfte eines Trägers
% =============================================

\section*{Übungsaufgabe -- Auflagerkräfte eines Trägers}

Ein an seinen beiden Enden gelagerter Träger der Länge~$L$ wird durch eine senkrecht wirkende Einzellast~$F$ beansprucht. Der Angriffspunkt dieser Last befindet sich im Abstand~$a$, gemessen vom linken Auflager~$A$; das gegenüberliegende, rechte Auflager wird mit~$B$ bezeichnet (siehe nachstehende Skizze).

\begin{center}
	\begin{tikzpicture}[>=Stealth, scale=1]
		% Balken
		\draw[line width=0.8pt, fill=gray!8] (0,0) rectangle (8,0.25);

		% Auflager A (links)
		\draw[fill=gray!25] (0,0) -- (-0.4,-0.65) -- (0.4,-0.65) -- cycle;
		\draw[line width=0.6pt] (-0.55,-0.65) -- (0.55,-0.65);
		\foreach \x in {-0.4,-0.2,...,0.45} {
			\draw[line width=0.4pt] (\x,-0.65) -- (\x-0.18,-0.88);
		}
		\node at (-0.85,-0.3) {$A$};

		% Auflager B (rechts)
		\draw[fill=gray!25] (8,0) -- (7.6,-0.65) -- (8.4,-0.65) -- cycle;
		\draw[line width=0.6pt] (7.45,-0.65) -- (8.55,-0.65);
		\foreach \x in {7.6,7.8,...,8.45} {
			\draw[line width=0.4pt] (\x,-0.65) -- (\x-0.18,-0.88);
		}
		\node at (8.85,-0.3) {$B$};

		% Einzellast F bei a = 3
		\draw[->, line width=1.1pt, red!70!black] (3,1.4) -- (3,0.25);
		\node[red!70!black] at (3,1.65) {$F$};

		% Maßlinie a
		\draw[<->, line width=0.5pt] (0,-1.25) -- (3,-1.25);
		\node[fill=white, inner sep=1pt] at (1.5,-1.25) {$a$};
		\draw[dotted] (3,0) -- (3,-1.35);

		% Maßlinie L
		\draw[<->, line width=0.5pt] (0,-1.85) -- (8,-1.85);
		\node[fill=white, inner sep=1pt] at (4,-1.85) {$L$};
		\draw[dotted] (0,-0.9) -- (0,-1.95);
		\draw[dotted] (8,-0.9) -- (8,-1.95);
	\end{tikzpicture}
\end{center}

Für einen derart belasteten Träger lassen sich die in den beiden Lagern auftretenden Auflagerkräfte~$F_A$ (am linken Lager) und~$F_B$ (am rechten Lager) aus dem Momentengleichgewicht herleiten und mit Hilfe der nachfolgenden Beziehungen bestimmen, welche für die Bearbeitung verwendet werden dürfen:

\[
	F_B = F \cdot \frac{a}{L}
	\qquad\qquad
	F_A = F \cdot \frac{L - a}{L}
\]

\begin{hinweis}[Physikalischer Hintergrund]
	Die Aufteilung der Last auf die beiden Lager richtet sich nach der Lage des Angriffspunktes: Je näher die Last am rechten Lager~$B$ angreift (großes~$a$), desto größer fällt der von~$B$ aufzunehmende Anteil~$F_B$ aus und desto kleiner wird~$F_A$. In jedem Fall entspricht die Summe beider Auflagerkräfte der aufgebrachten Last, es gilt also $F_A + F_B = F$.
\end{hinweis}

\medskip
In einem ersten Schritt werden über die Tastatur nacheinander \textbf{drei Zahlenwerte} eingelesen, und zwar für die Last~$F$, den Abstand~$a$ sowie die Trägerlänge~$L$.

\subsection*{Aufgabe}

\begin{enumerate}[label=\textbf{\arabic*.)}]
	\item Es soll geprüft werden, ob alle eingegebenen Zahlenwerte $\geq 0$ sind, ansonsten soll ein Fehlerhinweis ausgegeben werden.

	\item Die Trägerlänge $L$ soll größer als 0 sein, sonst solle eine Fehlermeldung angezeigt werden.

	\item Der Abstand $a$ soll als sinnvoller Angriffspunkt nicht außerhalb des Trägers liegen ($a \leq L$).

	\item Wenn keine Fehlermeldung nach 1, 2 und 3 vorliegt, sollen die Auflagerkräfte $F_A$ und $F_B$ berechnet und auf dem Bildschirm ausgegeben werden.
\end{enumerate}

\begin{hinweis}
	Für die Tastatureingabe kannst du die im Kurs verwendete Klasse \texttt{Tastatureingabe} mit der Methode \texttt{leseDouble()} nutzen. Alternativ ist auch ein \texttt{Scanner} erlaubt.
\end{hinweis}

\medskip
Vervollständige das folgende Programmgerüst:

\begin{lstlisting}
public class Auflagerkraft {

    public static void main(String[] args) {

        System.out.println("Geben Sie F ein:");
        double f = Tastatureingabe.leseDouble();

        System.out.println("Geben Sie a ein:");
        double a = Tastatureingabe.leseDouble();

        System.out.println("Geben Sie L ein:");
        double l = Tastatureingabe.leseDouble();
    }
}
\end{lstlisting}

\begin{beispiel}
	Eingabe: $F = 100$, $a = 2$, $L = 8$
	\begin{verbatim}
FA = 75.0
FB = 25.0
	\end{verbatim}
	(Die Last sitzt nahe am linken Lager, daher trägt $A$ den größeren Anteil.)
\end{beispiel}

\end{document}
